\documentclass[11pt]{article}
\usepackage{amsmath,amssymb,amsthm}
\usepackage{geometry}
\usepackage{braket}

\geometry{margin=1in}

\title{Project on Hamiltonian Simulation - Part II\\Theory Exercises}
\author{}
\date{}

\begin{document}
\maketitle

\section*{Exercise 1: An Algorithm for Trotterized Hamiltonian Simulation}

\subsection*{Exercise 1.1: Algorithm Sketch}

Recall the Trotter formula:
\begin{equation}
e^{-iHt} \approx \left( \prod_P e^{-ia_P P t/r} \right)^r
\end{equation}

Given the helper function \texttt{qc.pauli(P, theta)} implementing $e^{iP\theta}$, consider the following pseudocode:

\begin{verbatim}
procedure trotter_approx(H, r, t):
    # H is a list of (coefficient, pauli_string) pairs;
    # initialise quantum circuit with n qubits,
    # with n = Pauli string length
    n = len(H[0][1])
    qc = quantum_circuit(n)
    
    for step in range(r):
        for (a_P, P) in H:
            theta = -a_P * t / r
            qc.pauli(P, theta)
    
    return qc
\end{verbatim}

\textbf{Explanations:}
\begin{itemize}
    \item The outer range(r) ensures the power of r is applied
    \item The product over P, or application of all gates, is done in the inner loop using the helper function
    \item For theta, we need $e^{-ia_P P t/r}$, so we set $\theta = -a_P t/r$ so that $e^{iP\theta} = e^{iP(-a_P t/r)} = e^{-ia_P P t/r}$
\end{itemize}

\newpage
\subsection*{Exercise 1.2: Why Pauli Decomposition Format?}

There is a dimensionality problem when providing the full n qubit Hamiltonian.
\newline
Since the Hamiltonian is a matrix of dimensions $2^n \times 2^n$, the number of entries to store is $O(4^n)$. 
\newline
For large n (e.g. $n>30$), this gets exponentially large and becomes impossible to store!

\newline
\vspace{0.5cm}
On the other hand, for a Pauli decomposition of length $n$, and a Hamiltonian $H$ of \text{poly}(n) terms, we store $O(2*poly(n)*n) = O(poly(n))$ terms. This is much more manageable, as it is polynomial in length (and not exponential)!

\newline
\vspace{0.5cm}
The circuit then has $m$ $e^{-ia_P P t/r}$ exponentials constructed with $O(n)$ gates, repeated $r$ times, yielding complexity $O(mrn)$, polynomial in size. This is much better than the exponential complexity of the Hamiltonian derived earlier.

\newpage

\section*{Exercise 2: The Transverse-Field Ising Hamiltonian}

Recall the transverse-field Ising Hamiltonian formula:
\begin{equation}
H_{\text{Ising}} = -J \sum_{i=0}^{n-1} Z_i Z_{i+1(\text{mod }n)} - h \sum_i X_i
\end{equation}


\subsection*{2.1 Interaction Terms}

\subsubsection*{Part (a): proof that $Z_i Z_j$ is diagonal}

\textbf{Lemma:} The tensor product of diagonal matrices is diagonal.

\textbf{Proof:} Let $A = \text{diag}(a_1, \ldots, a_m)$ and $B = \text{diag}(b_1, \ldots, b_n)$ be diagonal matrices.

We write the Kronecker product $A \otimes B$ in block form:
\begin{equation}
A \otimes B = \begin{pmatrix}
A_{11} B & A_{12} B & \cdots & A_{1m} B \\
A_{21} B & A_{22} B & \cdots & A_{2m} B \\
\vdots & \vdots & \ddots & \vdots \\
A_{m1} B & A_{m2} B & \cdots & A_{mm} B
\end{pmatrix}
\end{equation}

Since $A$ is diagonal, $A_{ij} = 0$ for $i \neq j$, so all off-diagonal blocks vanish:
\begin{equation}
A \otimes B = \begin{pmatrix}
a_1 B & 0 & \cdots & 0 \\
0 & a_2 B & \cdots & 0 \\
\vdots & \vdots & \ddots & \vdots \\
0 & 0 & \cdots & a_m B
\end{pmatrix}
\end{equation}

Since $B$ is also diagonal, each block $a_i B$ is diagonal. Therefore, $A \otimes B$ is block-diagonal with diagonal blocks, making the entire matrix diagonal. $\square$

\textbf{Proof of part (a):} By the lemma, since I and Z are diagonal, and $Z_i, Z_j$ are tensor products of diagonal matrices, they are both diagonal in the computational basis.

\subsubsection*{Part (b): Action on computational basis states}

We compute how $e^{i(Jt/r) Z_i Z_{i+1}}$ acts on basis states $\ket{a}_i \ket{b}_{i+1}$ for $a, b \in \{0,1\}$.

First, recall that the Pauli $Z$ matrix acts as:
\begin{align}
Z\ket{0} &= \ket{0} \\
Z\ket{1} &= -\ket{1}
\end{align}

For the product $Z_i Z_{i+1}$:
\begin{align}
(Z_i Z_{i+1}) \ket{0}_i \ket{0}_{i+1} &= Z_i\ket{0} \otimes Z_{i+1}\ket{0} = \ket{0} \otimes \ket{0} = (+1) \ket{00} \\
(Z_i Z_{i+1}) \ket{0}_i \ket{1}_{i+1} &= Z_i\ket{0} \otimes Z_{i+1}\ket{1} = \ket{0} \otimes (-\ket{1}) = (-1) \ket{01} \\
(Z_i Z_{i+1}) \ket{1}_i \ket{0}_{i+1} &= Z_i\ket{1} \otimes Z_{i+1}\ket{0} = (-\ket{1}) \otimes \ket{0} = (-1) \ket{10} \\
(Z_i Z_{i+1}) \ket{1}_i \ket{1}_{i+1} &= Z_i\ket{1} \otimes Z_{i+1}\ket{1} = (-\ket{1}) \otimes (-\ket{1}) = (+1) \ket{11}
\end{align}

The eigenvalues are $\lambda_{ab} = (-1)^{a+b}$: $+1$ when $a = b$, and $-1$ when $a \neq b$.

From Part I Exercise 1.1, we proved that for a diagonal matrix $D = \text{diag}(\lambda_1, \ldots, \lambda_n)$:
\begin{equation}
e^D = \text{diag}(e^{\lambda_1}, \ldots, e^{\lambda_n})
\end{equation}

Since $Z_i Z_{i+1}$ is diagonal with entries equal to the eigenvalues computed above, we apply the exponential:
\begin{equation}
e^{i(Jt/r) Z_i Z_{i+1}} \ket{a}_i \ket{b}_{i+1} = e^{i(Jt/r) \lambda_{ab}} \ket{a}_i \ket{b}_{i+1} = e^{i(Jt/r)(-1)^{a+b}} \ket{a}_i \ket{b}_{i+1}
\end{equation}

Expanding, we get:
\begin{align}
e^{i(Jt/r) Z_i Z_{i+1}} \ket{0}\ket{0} &= e^{i(Jt/r)} \ket{0}\ket{0} \\
e^{i(Jt/r) Z_i Z_{i+1}} \ket{0}\ket{1} &= e^{-i(Jt/r)} \ket{0}\ket{1} \\
e^{i(Jt/r) Z_i Z_{i+1}} \ket{1}\ket{0} &= e^{-i(Jt/r)} \ket{1}\ket{0} \\
e^{i(Jt/r) Z_i Z_{i+1}} \ket{1}\ket{1} &= e^{i(Jt/r)} \ket{1}\ket{1}
\end{align}

\subsection*{2.2 Transverse Field Terms}

\subsubsection*{Part (a): Diagonalizing basis}

Recall the Pauli $X$ matrix has eigenvectors:
\begin{align}
\ket{+} &= \frac{1}{\sqrt{2}}(\ket{0} + \ket{1}), \quad \text{eigenvalue } +1 \\
\ket{-} &= \frac{1}{\sqrt{2}}(\ket{0} - \ket{1}), \quad \text{eigenvalue } -1
\end{align}

Applying the transverse field term $e^{i(ht/r) X_i}$ to each state:

\begin{align}
e^{i(ht/r) X_i} \ket{+} &= \sum_{k=0}^{\infty} \frac{(i(ht/r) X_i)^k}{k!} \ket{+} \\
&= \sum_{k=0}^{\infty} \frac{(i(ht/r))^k X_i^k}{k!} \ket{+} \\
&= \sum_{k=0}^{\infty} \frac{(i(ht/r))^k \cdot 1^k}{k!} \ket{+} \quad \text{(since } X_i\ket{+} = \ket{+}\text{)} \\
&= \left(\sum_{k=0}^{\infty} \frac{(i(ht/r))^k}{k!}\right) \ket{+} \\
&= e^{i(ht/r)} \ket{+}
\end{align}

Similarly:
\begin{align}
e^{i(ht/r) X_i} \ket{-} &= \sum_{k=0}^{\infty} \frac{(i(ht/r))^k X_i^k}{k!} \ket{-} \\
&= \sum_{k=0}^{\infty} \frac{(i(ht/r))^k \cdot (-1)^k}{k!} \ket{-} \quad \text{(since } X_i\ket{-} = -\ket{-}\text{)} \\
&= e^{-i(ht/r)} \ket{-}
\end{align}

Therefore $e^{i(ht/r)X_i}$ is diagonal in the $X$-basis $\{\ket{+}, \ket{-}\}$, with eigenvalues $e^{\pm i(ht/r)}$.
\subsubsection*{Part (b): Relation to rotation gate and bit flip}

Recall the rotation gate's definition:
\begin{equation}
R_X(\theta) = e^{-i\theta X/2}
\end{equation}

Compare with $e^{i(ht/r) X_i}$:
\begin{align}
e^{i(ht/r) X_i} &= e^{-i\theta X_i/2}
\end{align}

Matching exponents:
\begin{equation}
i(ht/r) = -i\theta/2 \quad \Rightarrow \quad \theta = -\frac{2ht}{r}
\end{equation}

Therefore:
\begin{equation}
e^{i(ht/r) X_i} = R_{X_i}\left(-\frac{2ht}{r}\right)
\end{equation}

\textbf{When the rotation gate acts as a bit flip}

The bit flip gate $X$ is equivalent to a rotation of (any odd multiple of) $\pi$:
\begin{equation}
X = R_X(\pi)
\end{equation}

For $e^{i(ht/r) X_i}$ to act as a bit flip:
\begin{equation}
-\frac{2ht}{r} = \pi + 2\pi k, \quad k \in \mathbb{Z}
\end{equation}

Solving for $t \geq 0$
\begin{equation}
t = -\frac{\pi r(1 + 2k)}{2h}
\end{equation}

For positive $t$, taking $k = -1, -2, -3, \ldots$:
\begin{equation}
t = \frac{\pi r(2m - 1)}{2h}, \quad m = 1, 2, 3, \ldots
\end{equation}

Therefore, the transverse field term acts as a bit flip when:
\begin{equation}
t \in \left\{ \frac{\pi r}{2h}, \frac{3\pi r}{2h}, \frac{5\pi r}{2h}, \ldots \right\}
\end{equation}


\end{document}